\begin{enumerate}
\item
\begin{enumerate}
    \item Au bout d'un tour, on obtient $1 + 5 = 1$. Au bout de deux tours, on obtient $1 + 10 = 1$. En ajoutant encore 2, il vient $1 + 12 = 3$. Visuellement, on a fait 2 tours d'horloge + 2 crans.
    
    \item Avec $n = 11$, au bout d'un tour, on a $1 + 11 = 1$. Si on avance encore d'un cran, cela donne $1 + 12 = 2$. Cette valeur de $n$ est la seule possible. En effet, si l'équation est vérifiée 12 vaut 1 additionné à un multiple de $n$, et donc $n$ divise 11, qui est premier, donc $n = 11$.
\end{enumerate}

\item \textbf{}
\begin{enumerate}
    \item C'est certain, Beethoven aurait été un compositeur modeste s'il avait été bien entendant.
    
    \item Un arrêté municipal interdit le ramassage des escargots. Une association demande de revenir sur cette décision. La mairie reste inflexible.
\end{enumerate}

\item \textbf{} Avec une demi-coupe, les cotes sont diminuées de moitié dans toutes les directions. Donc le volume est divisé par 8. On peut donc remplir 8 demi-coupes avec une coupe entière.

\item \textbf{} On superpose le verre penché et le verre droit.

Avec les notations, on a $AB = \dfrac{AC}{2}$, car $AC$ est le rayon du cercle, et car le verre est rempli à moitié.

Donc $\sin(x) = \dfrac{AB}{AC} = \dfrac{1}{2}$, donc l'angle maximum vaut $x = 30^\circ$.

\item \textbf{}
\begin{enumerate}
    \item Si $n$ est supérieur ou égal à 100 et si chacun paye déjà son menu à 40 euros, cela donne $40n$ euros qui est bien supérieur à 4000 euros : l'adjecteur est restauration et atteint et les convives n'auront pas besoin de payer davantage.
    
    Si $n$ est strictement inférieur à 100 et si chacun paye son menu à 40 euros, cela donne au total $40n$ euros qui est inférieur strictement à 4000 euros. Selon les conditions du restaurateur, les convives vont devoir payer un peu plus pour que la recette attaque bien 4000 euros. Ils payeront $4000/n$ euros par personne.
    
    Lorsque $n = 90$, on a vu que les convives vont devoir payer chacun $4000/90$ euros. C'est compris entre 4.44 euros et 4.45 euros : chacun payera donc 4.46 euros.
    
    \item On répartit la subvention sur les 90 convives. Selon le calcul précédent, chacun payera donc $4000/90 - 1000/90$ : cela donne 33.34 euros par personne.
    
    Lorsque 120 personnes viennent chacun payer $40 - 1000/120$ euros (chacun payer son menu -- sa part de subvention). Cela donne 31.67 euros.
    
    Si $n$ est supérieur ou égal à 100, le prix par convive est $40 - 1000/n$. Cette quantité est minimale pour $n = 100$ et vaut 30 (euros).
    
    Si $n$ est inférieur strictement à 100, le prix par convive est $4000/n - 1000/n$ c'est-à-dire $3000/n$. Cette quantité est minimale pour $n = 99$ et est supérieure à 30 (car $30 = 3000/100$).
    
    Conclusion : Le prix par convive est minimal lorsqu'il y a exactement 100 convives et le prix à payer par chacun est alors de 30 euros.
\end{enumerate}

\item \textbf{} Notons $R$ le rayon de la tarte. Et $x \in [0, R]$ l'endroit du bord où l'on va couper. L'aire de la portion totale de la vue vaut $\dfrac{\pi R^2}{4}$ et l'aire de la portion gauche de tarte vaut $\dfrac{x^2}{2}$. Ce dernier doit valoir la moitié de l'aire totale.

Donc $x = \dfrac{\sqrt{\pi}}{2} R$.

\item \textbf{} On note $S$ le sommet d'où partent les trois arêtes à angles droits, on note $A$, $B$ et $C$ les autres sommets ($C$ étant en arrière) et enfin, on note $a = SA$, $b = SB$ et $c = SC$. Appelons $d$ la hauteur issue de $B$ du triangle oblique $ABC$, et $H$ le pied de cette hauteur. L'aire au carré du triangle oblique vaut $AC^2  d^2 / 4$. Puis dans le triangle rectangle $HSB$, $d^2 = SH^2 + b^2$. Mais $(SH)$ est aussi perpendiculaire à $(AC)$. Donc l'aire du triangle de base, à savoir le triangle $ASC$, vaut à la fois $\dfrac{ac}{2}$ et $\dfrac{SH \cdot AC}{2}$. Il reste à tout recoller puisque $AC^2 = a^2 + c^2$. Dès lors, $SH^2(a^2 + c^2) = a^2c^2$. Et donc $d^2(a^2 + c^2) = SH^2(a^2 + c^2) + b^2(a^2 + c^2) = a^2c^2 + b^2c^2 + a^2b^2$.

\end{enumerate}
